\documentclass{article}
\usepackage[T1]{fontenc}
\usepackage{lmodern}
\usepackage{graphicx}
\usepackage{amsmath,amssymb}
\usepackage[a4paper,margin=2cm]{geometry}
\usepackage[absolute,overlay]{textpos}
\setlength{\TPHorizModule}{1mm}
\setlength{\TPVertModule}{1mm}
\usepackage[indonesian]{babel}
\usepackage{hyperref}
\setlength{\emergencystretch}{2em}
% unit: Halaman 6

\begin{document}

% === Halaman 6 (python/native) ===

6

• Contoh: misalkan pesan X = ‘AABBCBDB’


   n = 4  (yaitu huruf A, B, C, D)


   p(A) = 2/8, p(B) = 4/8


   p(C) = 1/8, p(D) = 1/8 


   H(x) = -\{2/8 2log(2/8) + 4/8 2log(4/8) + 1/8 2log(1/8) + 1/8 2log(1/8) \}

            = -\{1/4 2log(1/4) + 1/2 2log(1/2) + 1/8 2log(1/8) + 1/8 2log(1/8) \}

            = -\{(1/4) (- 2log(4) + (1/2) (- 2log(2) + (1/8) (- 2log(8)  + (/8) (- 2log(8) \}

            = -\{(1/4) (- 2) + (1/2) (-1) + (1/8) (-3) + (1/8)(-3) \}

            = - \{-1/2 – 1/2 – 3/8 – 3/8)

            = - (-1.75)

            = 1.75

    Entropi = 1,75 bit per simbol

\end{document}
